\documentclass[aspectratio=169]{beamer}
\usetheme{Madrid}
\usecolortheme{whale}
\usepackage[utf8]{inputenc}
\usepackage[spanish]{babel}
\usepackage{amsmath,amsfonts,amssymb}
\usepackage{graphicx}
\usepackage{booktabs}
\usepackage{listings}
\usepackage{xcolor}

\lstdefinestyle{python}{
  language=Python,
  basicstyle=\ttfamily\footnotesize,
  backgroundcolor=\color{gray!10},
  frame=single,
  rulecolor=\color{gray!50},
  breaklines=true,
  keywordstyle=\color{blue!80!black}\bfseries,
  commentstyle=\color{green!50!black}\itshape,
  stringstyle=\color{red!70!black},
  showstringspaces=false,
  numbers=none,
  tabsize=4
}

\title{Métodos de Análisis de Datos 1}
\subtitle{Clase 12: Visualización con Python}
\author{Depto.\ de Matemática -- UNS}
\institute{Universidad Nacional del Sur}
\date{1er Cuatrimestre 2026}

\AtBeginSection[]{
  \begin{frame}{Contenidos}
    \tableofcontents[currentsection]
  \end{frame}
}

\begin{document}

\begin{frame}
  \titlepage
\end{frame}

\begin{frame}{Contenidos}
  \tableofcontents
\end{frame}

%=============================================================
\section{Ecosistema de visualización en Python}
%=============================================================

\begin{frame}{Principales librerías}
  \begin{center}
    \begin{tabular}{lll}
      \toprule
      \textbf{Librería} & \textbf{Nivel} & \textbf{Fortaleza} \\
      \midrule
      \texttt{matplotlib} & Bajo nivel & Control total, base de todo \\
      \texttt{seaborn} & Alto nivel & Gráficos estadísticos elegantes \\
      \texttt{pandas} (.plot) & Alto nivel & Gráficos rápidos desde DataFrames \\
      \texttt{plotly} & Interactivo & Dashboards, web \\
      \texttt{bokeh} & Interactivo & Web, grandes datos \\
      \bottomrule
    \end{tabular}
  \end{center}
  \bigskip
  \begin{block}{Recomendación para el curso}
    Comenzar con \texttt{matplotlib} para entender la base; usar \texttt{seaborn} para gráficos estadísticos en el análisis exploratorio de datos (EDA).
  \end{block}
\end{frame}

%=============================================================
\section{matplotlib}
%=============================================================

\begin{frame}{Estructura de una figura en matplotlib}
  \textbf{Jerarquía de objetos:}
  \begin{itemize}
    \item \texttt{Figure}: contenedor principal (la ``hoja'').
    \item \texttt{Axes}: área de graficado (el ``gráfico''). Una figura puede tener varios \texttt{Axes}.
    \item Elementos: líneas, barras, puntos, textos, leyenda, ejes.
  \end{itemize}
  \bigskip
  \textbf{Dos interfaces:}
  \begin{itemize}
    \item \textbf{Orientada a objetos (OO)}: \texttt{fig, ax = plt.subplots()} — recomendada.
    \item \textbf{pyplot}: \texttt{plt.plot()}, \texttt{plt.show()} — rápida pero menos flexible.
  \end{itemize}
  \bigskip
  \begin{alertblock}{Buena práctica}
    Usar siempre la interfaz OO cuando el gráfico tiene más de un panel o requiere personalización.
  \end{alertblock}
\end{frame}

\begin{frame}[fragile]{Histograma y boxplot con matplotlib}
\begin{lstlisting}[style=python]
import matplotlib.pyplot as plt
import numpy as np

datos = np.random.normal(loc=50, scale=10, size=200)

fig, axes = plt.subplots(1, 2, figsize=(10, 4))

# Histograma
axes[0].hist(datos, bins=20, color='steelblue',
             edgecolor='white', density=True)
axes[0].set_title('Histograma')
axes[0].set_xlabel('Valor')
axes[0].set_ylabel('Densidad')

# Boxplot
axes[1].boxplot(datos, vert=True, patch_artist=True,
                boxprops=dict(facecolor='steelblue'))
axes[1].set_title('Boxplot')

plt.tight_layout()
plt.savefig('figura.png', dpi=150)
plt.show()
\end{lstlisting}
\end{frame}

\begin{frame}[fragile]{Gráfico de dispersión con matplotlib}
\begin{lstlisting}[style=python]
import matplotlib.pyplot as plt
import numpy as np

x = np.random.normal(0, 1, 100)
y = 2*x + np.random.normal(0, 0.5, 100)

fig, ax = plt.subplots(figsize=(6, 5))

ax.scatter(x, y, color='steelblue', alpha=0.6,
           edgecolors='white', s=60)

# Linea de tendencia
m, b = np.polyfit(x, y, 1)
ax.plot(sorted(x), [m*xi + b for xi in sorted(x)],
        color='red', linewidth=2, label='Tendencia')

ax.set_xlabel('Variable X')
ax.set_ylabel('Variable Y')
ax.set_title('Gráfico de dispersión')
ax.legend()
plt.tight_layout()
plt.show()
\end{lstlisting}
\end{frame}

%=============================================================
\section{seaborn}
%=============================================================

\begin{frame}{seaborn: introducción}
  \begin{block}{¿Qué es seaborn?}
    Librería de alto nivel construida sobre \texttt{matplotlib}, diseñada para crear gráficos estadísticos atractivos con poco código.
  \end{block}
  \bigskip
  \textbf{Características principales:}
  \begin{itemize}
    \item Trabaja nativamente con \texttt{pandas} DataFrames.
    \item Paletas de colores perceptualmente uniformes por defecto.
    \item Gráficos con separación por grupos (hue, col, row) muy sencilla.
    \item Incluye modelos estadísticos integrados (regresión, KDE).
  \end{itemize}
  \bigskip
  \textbf{Importación estándar:}
  \begin{itemize}
    \item \lstinline[style=python]|import seaborn as sns|
    \item \lstinline[style=python]|sns.set_theme(style='whitegrid')|
  \end{itemize}
\end{frame}

\begin{frame}[fragile]{Gráficos univariados con seaborn}
\begin{lstlisting}[style=python]
import seaborn as sns
import matplotlib.pyplot as plt

sns.set_theme(style='whitegrid')
tips = sns.load_dataset('tips')  # dataset de ejemplo

fig, axes = plt.subplots(1, 3, figsize=(12, 4))

# Histograma + KDE
sns.histplot(tips['total_bill'], kde=True,
             ax=axes[0], color='steelblue')

# Boxplot por grupo
sns.boxplot(x='day', y='total_bill', data=tips,
            ax=axes[1])

# Violin plot
sns.violinplot(x='day', y='total_bill', data=tips,
               ax=axes[2], inner='quartile')

plt.tight_layout()
plt.show()
\end{lstlisting}
\end{frame}

\begin{frame}[fragile]{Gráficos bivariados con seaborn}
\begin{lstlisting}[style=python]
import seaborn as sns
import matplotlib.pyplot as plt

sns.set_theme(style='whitegrid')
tips = sns.load_dataset('tips')

fig, axes = plt.subplots(1, 2, figsize=(10, 4))

# Scatter con regresion
sns.regplot(x='total_bill', y='tip', data=tips,
            ax=axes[0], scatter_kws={'alpha': 0.5})

# Heatmap de correlacion
corr = tips.select_dtypes('number').corr()
sns.heatmap(corr, annot=True, fmt='.2f',
            cmap='coolwarm', ax=axes[1])

plt.tight_layout()
plt.show()
\end{lstlisting}
\end{frame}

%=============================================================
\section{Buenas prácticas}
%=============================================================

\begin{frame}{Principios de visualización efectiva}
  \begin{enumerate}
    \item \textbf{Claridad}: cada gráfico debe comunicar un mensaje único y claro.
    \item \textbf{Títulos y etiquetas}: siempre incluir título, etiquetas de ejes y unidades.
    \item \textbf{Leyenda}: solo cuando hay más de una serie.
    \item \textbf{Colores}: usar paletas accesibles (daltónicas). Evitar el arcoíris (\textit{jet}).
    \item \textbf{Proporciones}: no truncar el eje $y$ arbitrariamente.
    \item \textbf{Grillas}: usar con moderación; no sobrecargar.
    \item \textbf{Resolución}: guardar en \texttt{.pdf} o \texttt{.svg} para publicaciones; \texttt{.png} con \texttt{dpi=150} o mayor.
  \end{enumerate}
\end{frame}

\begin{frame}{Tipos de gráficos y cuándo usarlos}
  \begin{center}
    \begin{tabular}{lll}
      \toprule
      \textbf{Gráfico} & \textbf{Datos} & \textbf{Para mostrar} \\
      \midrule
      Barras & Cualit. & Frecuencias, comparación categorías \\
      Histograma & Cuant. & Distribución \\
      Boxplot & Cuant. & Distribución + outliers + grupos \\
      Scatter & Cuant.+Cuant. & Asociación bivariada \\
      Línea & Tiempo+Cuant. & Tendencia temporal \\
      Heatmap & Matriz & Correlaciones, tablas de contingencia \\
      Violín & Cuant.+Grupos & Distribución por grupo \\
      \bottomrule
    \end{tabular}
  \end{center}
\end{frame}

%=============================================================
\section{Resumen}
%=============================================================

\begin{frame}{Resumen de la clase}
  \begin{block}{matplotlib}
    Base de la visualización en Python. Interfaz OO con \texttt{fig, ax = plt.subplots()}.
  \end{block}
  \bigskip
  \begin{block}{seaborn}
    Alto nivel sobre matplotlib. Gráficos estadísticos elegantes con DataFrames y separación por grupos (hue).
  \end{block}
  \bigskip
  \begin{block}{Próxima clase}
    Clase 13: Visualización con R — el sistema \texttt{ggplot2} basado en la Gramática de los Gráficos.
  \end{block}
\end{frame}

\end{document}
